
\begin{figure}[t]
\centering
\scriptsize
\resizebox{0.96\linewidth}{!}{%
\begin{tikzpicture}[
    >=Latex,
    source/.style={draw=blue!60!black, fill=blue!6, rounded corners, text width=2.7cm, minimum height=0.95cm, align=center},
    flow/.style={draw=orange!75!black, fill=orange!8, rounded corners, text width=2.9cm, minimum height=0.95cm, align=center},
    result/.style={draw=green!50!black, fill=green!8, rounded corners, text width=2.8cm, minimum height=0.95cm, align=center},
    aux/.style={draw=black!55, fill=gray!8, rounded corners, text width=2.55cm, minimum height=0.9cm, align=center}
]
\node[source] (hadm) at (-5.0,0.9) {$(\mathcal H_{\mathrm{adm}},\mathfrak A_{\mathrm{adm}})$\\admissible sector};
\node[flow] (wk) at (-1.6,0.9) {$W_k[J]$\\regulated generator};
\node[flow] (gk) at (1.8,0.9) {$\Gamma_k[\phi]$\\exact admissible flow};
\node[result] (gren) at (5.2,0.9) {$\Gamma_{\mathrm{TFPT}}^{\mathrm{ren}}$\\infrared 1PI action};

\node[aux] (alpha) at (-1.6,-1.25) {$\alpha_\ast$\\IR boundary value};
\node[aux] (shadow) at (1.8,-1.25) {one-loop shadow\\$\left(\frac{41}{10},-\frac{19}{6},-7\right)$};
\node[result] (obs) at (5.2,-1.25) {$\OphysTFPT$\\pole and response observables};
\node[result] (scat) at (8.5,-1.25) {$S_{\mathrm{massive}},\ \mathcal I_{\mathrm{dress}}$\\massive scattering\\and dressed interface};

\draw[->, thick] (hadm) -- (wk);
\draw[->, thick] (wk) -- (gk);
\draw[->, thick] (gk) -- (gren);
\draw[->, thick] (gren) -- (obs);
\draw[->, thick] (gren) -- (scat);
\draw[->, thick, dashed] (alpha) -- (gk);
\draw[->, thick, dashed] (gk) -- (shadow);
\end{tikzpicture}}
\caption{Exact admissible flow from the physical sector to observables and scattering. $\Padm$
selects states; $\Gamma_k$ carries running couplings and 1PI graphs.}
\label{fig:admissible-rg-flow}
\end{figure}

\begin{corollary}[No double counting after the seed-shadow split]
\label{cor:no-double-counting}
Rows sharing the same basis package in \Cref{def:minimal-observable-basis} define one underlying
operational direction, not several independent confirmations or failures. In particular, the former
quartet
\[
(\beta,\Omega_b,\lambda_C,\sin^2\theta_{13})
\]
no longer constitutes one seed package. Instead
\[
\beta\in R_{\mathrm{CS}},
\qquad
\Omega_b\in C_{\mathrm{cos}},
\qquad
\lambda_C\in F_{\mathrm{fl}},
\qquad
\sin^2\theta_{13}\in N_\nu,
\]
while the retained seed $\phiz$ belongs only to the UV kernel package $K_{\mathrm{UV}}$.
\end{corollary}

\begin{proof}
This is the special case of \textbf{(O2)} for the rewritten basis packages of
\Cref{def:minimal-observable-basis}.
\end{proof}

\begin{corollary}[Independent response-sector birefringence target]
\label{cor:response-birefringence-target}
On the minimal branch the determinant-response package $R_{\mathrm{CS}}$ fixes
\[
\beta_{\mathrm{rad}}=\frac{\phi_{\mathrm{ret}}^0}{4\pi},
\]
and hence the dimensionless birefringence readout itself. Under the declared comparison interface it
is represented by
\[
\beta
=
\frac{180}{\pi}\,\beta_{\mathrm{rad}}
=
0.2424350309009295\ldots^\circ.
\]
This target factors through $R_{\mathrm{CS}}$ and not through the flavor package
$F_{\mathrm{fl}}$ or the cosmology package $C_{\mathrm{cos}}$.
\end{corollary}

\begin{proof}
By the exact seed relation of the determinant-response channel,
\[
\beta_{\mathrm{rad}}=\frac{\phi_{\mathrm{ret}}^0}{4\pi}.
\]
The theorem-level quantity is therefore the dimensionless angle $\beta_{\mathrm{rad}}$. Converting
to degrees gives one scheme-facing representative,
\[
\beta=\frac{180}{\pi}\beta_{\mathrm{rad}}
=
\frac{180}{\pi}\frac{\phi_{\mathrm{ret}}^0}{4\pi}
=
0.2424350309009295\ldots^\circ.
\]
By \TFPTxref{cor:no-double-counting}, $\beta$ belongs to $R_{\mathrm{CS}}$, whereas
$\lambda_C\in F_{\mathrm{fl}}$ and $\Omega_b\in C_{\mathrm{cos}}$. Hence this is an operational
direction distinct from the flavor and cosmology packages.
\end{proof}

\begin{corollary}[Independent cosmology-sector axion target]
\label{cor:cosmology-axion-target}
On the closed cosmology branch the package $C_{\mathrm{cos}}$ fixes the axion readout as one
dimensionless cosmology package, equivalently as the boundary-normalized row
\[
\left(
\frac{f_a}{\lambda_\Sigma},
\frac{m_a}{\lambda_\Sigma},
\frac{\nu_a}{\lambda_\Sigma},
\lambda_\Sigma\left|g_{a\gamma\gamma}^{\mathrm{phys}}\right|
\right),
\]
This target factors through $C_{\mathrm{cos}}$ and not through the flavor package
$F_{\mathrm{fl}}$ or the determinant-response package $R_{\mathrm{CS}}$.
\end{corollary}

\begin{proof}
The closed cosmology branch fixes the determinant-line / seam-transfer / scalaron saddle and thus
the axion row as a single component of $C_{\mathrm{cos}}$. Its unitful representatives belong only
to the declared comparison interface and are recorded in the appendix benchmark layer, together with
the practical haloscope window and coupling representative extracted from the same cosmology
package.
By \TFPTxref{cor:no-double-counting}, this row belongs to $C_{\mathrm{cos}}$, whereas $\beta$ belongs
to $R_{\mathrm{CS}}$ and $\lambda_C$ belongs to $F_{\mathrm{fl}}$. Hence the axion target is an
operational direction distinct from response and flavor.
\end{proof}

\begin{remark}[Why these two targets count as independent]
The pair
\[
\beta\in R_{\mathrm{CS}},
\qquad
(\nu_a,m_a,g_{a\gamma\gamma}^{\mathrm{phys}})\in C_{\mathrm{cos}},
\]
samples two different basis packages of $\Bind(\Tstar)$. Therefore they are not two rewritings of
the same seed constraint. This is the correct sharp-prediction surface for the present version.
\end{remark}


% ---- Reflection positivity, OS reconstruction, and scattering appendices: source lines 13934-14109 ----
\section{Relative reflection positivity with fermions}
\label{app:os-fermions}

\subsection{Normalized quotient form}

The relative Euclidean theory used in the main closure section is not interpreted as a difference of measures. Instead it is organized as the normalized quotient
\[
Z_{\mathrm{rel}}[J,\eta,\bar\eta]
=
\frac{Z_{\mathrm{adm}}[J,\eta,\bar\eta]}{Z_{\mathrm{ref}}[0,0,0]}.
\]
Because the denominator is source independent on the admissible sector, the reference theory acts only as a fixed normalizer and does not interfere with reflection positivity.

\subsection{Bosonic and fermionic positivity on \texorpdfstring{$\Padm$}{Padm}}

The commutation relations
\[
[\Theta,P_\Theta]=0,
\qquad
\Theta\Padm=\Padm\Theta
\]
ensure that reflection preserves the admissible subspace. The bosonic part then follows the standard Markov / reflection-stable Osterwalder--Schrader argument on admissible observables. For fermions, the admissible CAR kernel may be written as a reflected quadratic form; equivalently one may express the same condition as positivity of the associated Pfaffian structure on the admissible sector. Hence for every admissible observable $O$ supported in $\tau>0$,
\[
\langle \Theta O,O\rangle_{\mathrm{rel}}
=
\frac{1}{Z_{\mathrm{ref}}[0,0,0]}
\langle \Theta O,O\rangle_{\mathrm{adm}}
\ge 0.
\]
This appendix section is the mechanism behind \TFPTxref{lem:admissible-reflection-positivity,thm:reflection-positivity-admissible,thm:lorentzian-reconstruction}; it is recorded here so that the main-text reconstruction theorem is not supported only by a one-line axiom call.

\section{Admissible OS reconstruction and asymptotic scattering}
\label{app:os-scattering}

This appendix upgrades the admissible Schwinger, reconstruction, and scattering block from a named
standard interface to a manuscript-level proof package. The point is not to replace standard
machinery by ad hoc notation, but to verify explicitly that the admissible branch satisfies the
relevant hypotheses on $\Padm$.

\subsection{Tempered admissible Schwinger families}

\begin{lemma}[Temperedness and moment bounds for admissible Schwinger functions]
\label{lem:tempered-schwinger-bounds}
On the Euclidean continuation of $\mathcal U_{\mathrm{scat}}$, each truncated admissible Schwinger
distribution $S_n^T$ defines a tempered distribution. Moreover, for every compact $K$ and every
source monomial of total degree $n$, its moments obey polynomial bounds determined by the same
uniform heat-kernel estimates and finite local moments of the admissible constructive measure.
\end{lemma}

\begin{proof}
By \TFPTxref{thm:well-posed-primitive-dynamics}, the admissible Euclidean generator is realized by a
self-adjoint elliptic operator on every blocked volume. Together with the constructive measure of
\TFPTxref{thm:constructive-geometric-measure}, this yields uniform heat-kernel bounds and finite moments
for local source insertions. Differentiating $W_{\mathrm{rel}}$ with respect to the local sources
therefore produces distributions whose derivatives are controlled by polynomial combinations of the
same bounds. Passing to the admissible projective limit preserves these polynomial estimates, so the
resulting Schwinger family is tempered.
\end{proof}

\begin{lemma}[Graded Euclidean covariance and symmetry on \texorpdfstring{$\Padm$}{Padm}]
\label{lem:graded-euclidean-covariance-padm}
On the Euclidean continuation of $\mathcal U_{\mathrm{scat}}$, the admissible Schwinger family is
graded Euclidean covariant, graded symmetric under exchange of external legs, and continuous in the
Schwartz topology.
\end{lemma}

\begin{proof}
The Dirac-type principal symbol of the geometric anchor and the low-curvature reduction imply local
Euclidean covariance of the admissible source functional on $\mathcal U_{\mathrm{scat}}$. Bosonic
source derivatives commute, while fermionic derivatives anticommute, so the mixed Schwinger family
is graded symmetric with the usual CAR signs. Continuity follows from the tempered moment bounds of
\Cref{lem:tempered-schwinger-bounds}, which control the action of $S_n^T$ on Schwartz test
functions.
\end{proof}

\subsection{OS and CAR reconstruction on \texorpdfstring{$\Padm$}{Padm}}

\begin{theorem}[Tailored OS/CAR reconstruction on \texorpdfstring{$\Padm$}{Padm}]
\label{thm:tailored-os-car-reconstruction}
Let $\{S_n^T\}_{n\ge 1}$ be the admissible Schwinger family on the Euclidean continuation of
$\mathcal U_{\mathrm{scat}}$. If reflection positivity, temperedness, graded Euclidean covariance,
graded symmetry, and clustering hold on $\Padm$, then there exist:
\begin{enumerate}[label=(\roman*), leftmargin=1.8em, topsep=3pt, itemsep=2pt]
    \item a Hilbert space $\mathcal H_{\mathrm{OS}}$ with vacuum vector $\Omega$,
    \item a strongly continuous positive-energy semigroup $T_t=e^{-tH_{\mathrm{adm}}}$ with
    self-adjoint generator $H_{\mathrm{adm}}\ge 0$,
    \item a reconstructed graded local field net on $\Padm$ whose Euclidean boundary values are the
    original admissible Schwinger functions.
\end{enumerate}
Equivalently, the admissible Euclidean family reconstructs a Wightman/CAR theory on the admissible
sector.
\end{theorem}

\begin{proof}
Define the OS sesquilinear form by
\[
(F,G)_{\mathrm{OS}}:=\langle \Theta F,G\rangle_{\mathrm{rel}}
\]
for positive-time admissible observables. Reflection positivity from
\TFPTxref{thm:reflection-positivity-admissible} and the fermionic quotient-form analysis of
\Cref{app:os-fermions} show that this form is positive semidefinite on the full graded admissible
algebra. Quotienting by the null space and completing gives $\mathcal H_{\mathrm{OS}}$.

Positive Euclidean time translation preserves the null space and acts by contractions, hence defines
a strongly continuous contraction semigroup $T_t$ on $\mathcal H_{\mathrm{OS}}$. The dense domain of
positive-time polynomial observables gives the symmetric generator domain. By Hille--Yosida, the
closure of this generator is self-adjoint and nonnegative, so $T_t=e^{-tH_{\mathrm{adm}}}$ for a
self-adjoint $H_{\mathrm{adm}}\ge 0$.

The temperedness, covariance, graded symmetry, and clustering established in
\Cref{lem:tempered-schwinger-bounds,lem:graded-euclidean-covariance-padm} are precisely the
remaining OS/CAR reconstruction hypotheses. Therefore analytic continuation yields the reconstructed
Wightman/CAR family with vacuum $\Omega$. Because $\Theta$, $P_\Theta$, and $\Padm$ commute, the
entire reconstruction stays on the admissible sector $\Padm$.
\end{proof}

\subsection{Stable massive Haag--Ruelle and LSZ closure}

\begin{lemma}[One-particle spectral projectors from isolated simple poles]
\label{lem:one-particle-projectors-from-poles}
Let $\widetilde G_a^{(2)}(p)$ have an isolated simple pole with positive residue at
$p^2=-m_a^2<0$. Then the corresponding one-particle contribution defines a positive spectral
projector on the admissible Hilbert space.
\end{lemma}

\begin{proof}
The isolated simple pole separates a one-particle mass shell from the multiparticle continuum. By
positivity of the residue, the contour integral of the resolvent around that pole defines a positive
projector on $\mathcal H_{\mathrm{adm}}$. Its range is the one-particle subspace associated with the
stable species $a$.
\end{proof}

\begin{theorem}[Haag--Ruelle scattering on the stable massive admissible sector]
\label{thm:haag-ruelle-admissible-massive}
For every stable massive admissible species with isolated simple pole below threshold, the
Haag--Ruelle asymptotic fields exist on $\mathcal H_{\mathrm{adm}}$ and generate wave operators on
the stable massive scattering sector.
\end{theorem}

\begin{proof}
By \TFPTxref{thm:tailored-os-car-reconstruction,thm:admissible-local-minkowski-net}, the admissible
branch carries a local relativistic net with positive Hamiltonian. By
\TFPTxref{thm:nonperturbative-admissible-qft}, connected correlators cluster exponentially on the
massive sector. The pole assumption and \Cref{lem:one-particle-projectors-from-poles} identify the
one-particle subspace and give positive single-particle spectral projectors. The standard
Haag--Ruelle estimates therefore apply to admissible local operators, yielding the in/out fields and
the wave operators on the stable massive sector.
\end{proof}

\begin{corollary}[LSZ reduction on the stable massive admissible sector]
\label{cor:lsz-admissible-massive}
On the stable massive admissible sector, matrix elements of the scattering operator are obtained by
amputating the connected correlators with the corresponding residue factors.
\end{corollary}

\begin{proof}
By \Cref{thm:haag-ruelle-admissible-massive}, the asymptotic fields and wave operators exist on the
stable massive sector. Applying the usual LSZ reduction to those asymptotic fields and the
admissible local net yields the displayed amputated-correlator formula.
\end{proof}

\subsection{Dressed massless interface in the low-curvature region}

The massless sector is not treated by bare Fock asymptotics in the present manuscript. The relevant
object is the infrared-dressed interface package
\[
\mathcal I_{\mathrm{dress}}
:=
\bigl(\Omega_+^{\mathrm{dress}},\Omega_-^{\mathrm{dress}},S_{\mathrm{dress}}\bigr),
\qquad
S_{\mathrm{dress}}:=(\Omega_+^{\mathrm{dress}})^\ast\Omega_-^{\mathrm{dress}},
\]
constructed from the soft charges of the admissible local net on the low-curvature region
$\mathcal U_{\mathrm{scat}}$. This is exactly the interface package used in the main-text theorem
\TFPTxref{thm:dressed-massless-interface}.


% ---- Yukawa kernels and positivity lemmas: source lines 14191-14244 ----
\section{Yukawa kernels and positivity lemmas}
\label{app:yukawa}

This appendix is the proof motor for the transport positivity, determinant-phase suppression, and
bounded-residual statements used in the main text. The transport program is based on the regularized kernel
\[
\mathcal Y_f^{(\varepsilon)}=(D_f^\dagger D_f+\varepsilon^2)^{-1}.
\]
Three remarks are important.
\begin{enumerate}[leftmargin=1.5em, topsep=3pt, itemsep=3pt]
    \item The regulator $\varepsilon$ is not optional: without it the kernel becomes ill-defined exactly where the admissibility question is most delicate.
    \item Positivity of the kernel belongs to the admissibility-closure layer rather than to the hard carrier kernel. The main text closes it through the finite hexagon gap and the lift to the full admissibility-projected operator; this appendix records the transport-side mechanism behind that closure.
    \item Residual prefactors $\Lambda_{f,j}$ are bounded holonomy data inside the transport theorem. They must remain visible and must not be rhetorically hidden as if the mass map were diagonal without residual transport structure.
\end{enumerate}

\begin{remark}[Positive-kernel determinant route to strong CP]
If the physical mass map is written in the polar form
\[
M_f=U_{f,L}H_fU_{f,R}^\dagger,
\qquad
H_f>0,
\qquad
U_{f,L},U_{f,R}\in SU(3)_F,
\]
with $H_f$ induced by the positive transport kernel $\mathcal Y_f^{(\varepsilon)}$, then the induced determinant phase vanishes:
\[
\arg\det M_f=0.
\]
At the same time weak CP may survive through
\[
V_{\mathrm{CKM}}=U_{u,L}^\dagger U_{d,L}.
\]
Combined with the sheet-CP symmetry protection of $\theta_{\mathrm{sheet}}$ in the relative action,
this is the algebraic core of the determinant-phase theorem stated in the main text and of the
strong-CP closure theorem recorded there. The role of the appendix remark is to display the
polar-decomposition mechanics behind that route, not to reopen its status.
\end{remark}

\subsection{Optional strong-coupling sharpening}

As an appendix-level dynamical sharpening of the hadronic admissibility theorem, one may further require the
admissible center-neutral Wilson functional to obey strong-coupling dominance. On that optional
route the relative Wilson loop satisfies
\[
\langle W(C)\rangle_{\mathrm{rel}}
\le
\exp[-\sigmaQCD A_{\min}(C)],
\qquad
\sigmaQCD=\cthree^2\chi_{\QCD}^2,
\]
and the physical Hilbert space contains no color-nonsinglet asymptotic states. This statement is
kept in the appendix so that the main theorem surface does not depend on an extra strong-coupling
sharpening beyond the closed admissible hadronic sector.

